\documentclass[conference]{IEEEtran}
\IEEEoverridecommandlockouts
% The preceding line is only needed to identify funding in the first footnote. If that is unneeded, please comment it out.
\usepackage{cite}
\usepackage{amsmath,amssymb,amsfonts}
\usepackage{algorithmic}
\usepackage{graphicx}
\usepackage{textcomp}
\usepackage{xcolor}
\usepackage{booktabs}

\def\BibTeX{{\rm B\kern-.05em{\sc i\kern-.025em b}\kern-.08em
    T\kern-.1667em\lower.7ex\hbox{E}\kern-.125emX}}
\begin{document}


\title{Short and Descriptive Title}
% write your names here
\author{\IEEEauthorblockN{Abc Xyz, Efg Pqr
% Change your group ID here and Option type
\IEEEauthorblockA{\textit{SYED 522 -- Group 2, Option A} \\
\textit{Systems Design Department}\\
\textit{University of Waterloo, Canada}\\
\texttt{{\{abc.xyz, efg.pqr}\}@uwaterloo}}
}}

\maketitle

%%%%%%%%%%%%%%%%%%%%%%%%%%%%%%%%%%%%%%%%%%%%%%%%%%%%%%%%%%%%%%%%%%%%%%%%%%%%%%%%%%%%%%%%%%%%%%%%%%%%%%%%%%%%%%%%%%%%%%%
\begin{abstract}
Abstract should be around ~300 words max. Do not cite any references here. Abstract and conclusion together are 10\% of the report (5\% each). Even though it is 5\% mark-wise but it is a really important section to grab attention of your readers. So pay attention while writing the abstract. Readers usually continue reading the paper based on the information provided in the abstract.

\end{abstract}

\begin{IEEEkeywords}
% 2-5 keywords don't use generics ones like Machine-Learning (be more specific)
Neural Networks, SVM, Linear Regression 
\end{IEEEkeywords}


% Sections
%%%%%%%%%%%%%%%%%%%%%%%%%%%%%%%%%%%%%%%%%%%%%%%%%%%%%%%%%%%%%%%%%%%%%%%%%%%%%%%%%%%%%%%%%%%%%%%%%%%%%%%%%%%%%%%%%%%%%%%
\section{Introduction} \label{sec:intro}
What is the topic? Why is this problem relevant? list the existing limitations. cite papers (only scientific papers - websites will be penalized) to back your claim. The Introduction section points out the gap in knowledge that the rest of the paper will fill, or the section in which you define and claim your territory within the broad area of research. The last paragraph of introduction generally describes how the paper is organized. \\

\noindent \textbf{Length:} Max 2 columns (1 page)\\
\noindent \textbf{Marks:} 10\%

%%%%%%%%%%%%%%%%%%%%%%%%%%%%%%%%%%%%%%%%%%%%%%%%%%%%%%%%%%%%%%%%%%%%%%%%%%%%%%%%%%%%%%%%%%%%%%%%%%%%%%%%%%%%%%%%%%%%%%%
\section{Background}\label{sec:background}
Background section contains the literature review of the existing approaches for solving the proposed problem. It may also describe any specific method that you're using (lets say SVM). Basically, you want readers to know these topics and the existing literature before describing them your method/study in details.

\subsection{SVM}
Support Vector Machine is a supervised machine learning method invested by Vladimir Vapnik. And more details and citations.

\subsection{ANN}
This is another subsection about ANNs. Backpropagation enables training. Add a peer-reviewed citation here. \\

If you're using the Option 1 as your project, this will be the major section for your report. Make sure, you divide this section into multiple sub-sections for the better organization. You should aim to have peer-reviewed citations (25-30 citations). Check the final project paper instructions for more details. \\

\noindent \textbf{Length:} 2 columns (Option 2, 3), 4-5 columns  (Option 1)\\
\noindent \textbf{Marks:} 20\% (Option 2, 3), 40\% (Option 1)

%%%%%%%%%%%%%%%%%%%%%%%%%%%%%%%%%%%%%%%%%%%%%%%%%%%%%%%%%%%%%%%%%%%%%%%%%%%%%%%%%%%%%%%%%%%%%%%%%%%%%%%%%%%%%%%%%%%%%%%
\section{Method} \label{sec:method}
If you're using Option 2 or 3 then describe your method here. If you are using Option 1 you should call this "Main State-of-the-Art". You may use figures and diagrams. Prefer high-quality/high-res images. For example, check the UW campus Figure~\ref{fig:uw-campus}.Check the final project paper instructions for more details. Check the final project paper instructions for more details.\\

\begin{figure}[tbh]
\centering
\IfFileExists{Figures/uw-campus.pdf}{
    \includegraphics[width=0.4\textwidth]{Figures/uw-campus}
}{
    \fbox{\parbox[c][1in][c]{0.4\textwidth}{\centering Figure placeholder\\Replace with a project figure.}}
}
\caption{UW Campus. Every caption should have a one liner summary followed by the detailed description. Every figure should be refereed within the text. \label{fig:uw-campus}}
\end{figure}


\noindent \textbf{Length:} 2-3 columns (Option 2, 3), 2 columns (Option 1) \\
\noindent \textbf{Marks:} 20\% (Option 2, 3), 10\% (Option 1)

%%%%%%%%%%%%%%%%%%%%%%%%%%%%%%%%%%%%%%%%%%%%%%%%%%%%%%%%%%%%%%%%%%%%%%%%%%%%%%%%%%%%%%%%%%%%%%%%%%%%%%%%%%%%%%%%%%%%%%%
\section{Results \& Analysis} \label{sec:results}
If you're using Option 1 then the section headline becomes ``Analysis'' only. Here you discuss the datasets used, experiments details, and the results. You may divide this section into multiple sub-sections as desired. This is another major section of the report. Be comprehensive about how you conducted the experiments, datasets, and the results and discussions. You should use graphs, plots, tables to present the results. For example, some results are presented in Table~\ref{tab:results}. Check the final project paper instructions for more details.\\

\begin{table}[bt]
\caption{Result on test data. Table captions in the template is at the top (this is different from figures). The caption formatting is the same way, one liner summary followed by the detailed description. Always refer you table inside the text.\label{tab:results}}
\centering
 \begin{tabular}{c c c c} 
 \toprule
 Col1 & Col2 & Col2 & Col3 \\ [0.5ex] 
 \midrule
 1 & 6 & 87837 & 787 \\ 
 2 & 7 & 78 & 5415 \\
 3 & 545 & 778 & 7507 \\
 4 & 545 & 18744 & 7560 \\
 5 & 88 & 788 & 6344 \\ [1ex] 
 \bottomrule
 \end{tabular}
\end{table}

\noindent \textbf{Length:} 4 columns (Option 2, 3), 2-3 columns (Option 1)\\
\noindent \textbf{Marks:} 35\% (Option 2, 3), 20\% (Option 1)

%%%%%%%%%%%%%%%%%%%%%%%%%%%%%%%%%%%%%%%%%%%%%%%%%%%%%%%%%%%%%%%%%%%%%%%%%%%%%%%%%%%%%%%%%%%%%%%%%%%%%%%%%%%%%%%%%%%%%%%
\section{Conclusions} \label{section:conclusion}
Summarize your report with concluding remarks. Some future studies or recommendations are strongly recommended.\\

\noindent \textbf{Length:} 1 column (1/2 page)\\
\noindent \textbf{Marks:} 10\% (Conclusion + Abstract together)

%%%%%%%%%%%%%%%%%%%%%%%%%%%%%%%%%%%%%%%%%%%%%%%%%%%%%%%%%%%%%%%%%%%%%%%%%%%%%%%%%%%%%%%%%%%%%%%%%%%%%%%%%%%%%%%%%%%%%%%
\section{Authors Contributions \& Remarks} 
No marks for this section. Feel free to skip this section. But if you want to tell us how each of the group members have contributed then add this section to your report. You may also add any other remarks that we should consider while marking your report. Keep it short and this won't be counted towards the final length of the report.

%%%%%%%%%%%%%%%%%%%%%%%%%%%%%%%%%%%%%%%%%%%%%%%%%%%%%%%%%%%%%%%%%%%%%%%%%%%%%%%%%%%%%%%%%%%%%%%%%%%%%%%%%%%%%%%%%%%%%%%
\section{references} \label{section:ref}
References here. Include only sources cited in the paper. \emph{This section by itself can be removed since it will be generated automatically.} They do \textbf{not} count towards your 6 pages limit.\\

\noindent \textbf{Length:} 2 columns (Option 2, 3), 3 columns (Option 1)\\
\noindent \textbf{Marks:} 5\% (Option 2, 3), 10\% (Option 1)

% Uncomment after adding the project's bibliography file and citations.
% \bibliographystyle{unsrt}
% \bibliography{references}


\end{document}
